\section{Extension and Delegation Inheritance}

\begin{frame}{Two Inheritance Styles}
	\begin{itemize}
		\item Odoo model inheritance has two major styles:
		\begin{enumerate}
			\item \textbf{Extension inheritance} with \texttt{\_inherit}
			\item \textbf{Delegation inheritance} with \texttt{\_inherits}
		\end{enumerate}
		\item They solve different problems.
		\item Understanding both is essential for clean module design.
	\end{itemize}
\end{frame}

\begin{frame}[fragile]{Extension Inheritance (\texttt{\_inherit})}
	\begin{itemize}
		\item Extends an existing model in place.
		\item Adds fields/methods to the same model/table.
\begin{block}{Method Syntax}
\begin{lstlisting}
class Student(models.Model):
	_inherit = 'student.student'

	blood_group = fields.Char()
\end{lstlisting}
\end{block}
		\item No new model name is required when you only extend.
	\end{itemize}
\end{frame}

\begin{frame}[fragile]{Classical Extension Patterns}
\begin{lstlisting}
# 1) Extend existing model
class ResPartner(models.Model):
	_inherit = 'res.partner'
	is_teacher = fields.Boolean()

# 2) Prototype inheritance: new model based on another
class StudentAlumni(models.Model):
	_name = 'student.alumni'
	_inherit = 'student.student'
\end{lstlisting}
	\begin{itemize}
		\item \texttt{\_inherit} alone extends the original model.
		\item \texttt{\_name} + \texttt{\_inherit} creates a new model copying the parent definition.
	\end{itemize}
\end{frame}

\begin{frame}[fragile]{Delegation Inheritance (\texttt{\_inherits})}
	\begin{itemize}
		\item Creates a new model that \textbf{delegates} fields to a linked parent record.
		\item Classic example: \texttt{res.users} delegates to \texttt{res.partner}.
\begin{block}{Method Syntax}
\begin{lstlisting}
class StudentProfile(models.Model):
	_name = 'student.profile'
	_inherits = {'res.partner': 'partner_id'}

	partner_id = fields.Many2one(
		'res.partner', required=True, ondelete='cascade')
	admission_no = fields.Char()
\end{lstlisting}
\end{block}
	\end{itemize}
\end{frame}

\begin{frame}[fragile]{How Delegation Works}
\begin{lstlisting}
profile = self.env['student.profile'].create({
	'name': 'Ahmed Muhumed',      # goes to res.partner
	'email': 'ahmed@example.com',  # goes to res.partner
	'admission_no': 'A-001',       # local field
})
print(profile.name)   # delegated field
print(profile.partner_id)
\end{lstlisting}
	\begin{itemize}
		\item Delegated fields are readable/writable through the child model.
		\item Data is stored in the parent table.
	\end{itemize}
\end{frame}

\begin{frame}[fragile]{Extension vs Delegation}
	\begin{itemize}
		\item Use \textbf{extension} when you want to customize an existing model.
		\item Use \textbf{delegation} when you want a new model that reuses another model's data compositionally.
	\end{itemize}
\begin{lstlisting}
_inherit = 'student.student'              # extend
_inherits = {'res.partner': 'partner_id'} # delegate
\end{lstlisting}
	\begin{itemize}
		\item Do not confuse \texttt{\_inherit} with \texttt{\_inherits}.
	\end{itemize}
\end{frame}

\begin{frame}[fragile]{Inheritance Best Practices}
	\begin{itemize}
		\item Prefer extension for small feature additions.
		\item Call \texttt{super()} when overriding methods.
		\item Keep delegated \texttt{Many2one} required with a clear \texttt{ondelete}.
		\item Avoid deep inheritance chains that are hard to debug.
	\end{itemize}
\end{frame}
